% answers_short.tex —— 2011 年·江苏卷 · 简短答案速查
% 由 tools/make_short_answers.py 自 answers.tex 的 \gksolhead/\gkanswer 提取，
% 并与 2011/2011_数学真题_江苏_答案要点.md 逐题比对后生成，勿手改（改源 MD 后重跑）。

\begin{gknotice}
\textbf{使用说明}：本节只给最终答案便于快速核对，完整推导见解析版。答案为编写者独立推导，\textbf{非官方标准答案}。
\end{gknotice}

\gkpart{填空题}
\begin{itemize}[leftmargin=2.6em]
  \item[\textbf{1.}] $\{-1,2\}$。
  \item[\textbf{2.}] $\left(-\dfrac12,+\infty\right)$。
  \item[\textbf{3.}] $1$。
  \item[\textbf{4.}] $3$。
  \item[\textbf{5.}] $\dfrac13$。
  \item[\textbf{6.}] $\dfrac{16}5$。
  \item[\textbf{7.}] $\dfrac49$。
  \item[\textbf{8.}] $4$。
  \item[\textbf{9.}] $\dfrac{\sqrt6}2$。
  \item[\textbf{10.}] $\dfrac54$。
  \item[\textbf{11.}] $-\dfrac34$。
  \item[\textbf{12.}] $\dfrac{\mathrm e+\mathrm e^{-1}}2$。
  \item[\textbf{13.}] $\sqrt[3]{3}$。
  \item[\textbf{14.}] $[1,2+\sqrt2]$。
\end{itemize}

\gkpart{解答题}
\begin{itemize}[leftmargin=2.6em]
  \item[\textbf{15.}] (1) $\dfrac\pi3$；(2) $\dfrac13$
  \item[\textbf{16.}] (1) 证明见解析；(2) 证明见解析
  \item[\textbf{17.}] (1) $x=15$；(2) $x=10$，高与底面边长之比为 $\dfrac14$
  \item[\textbf{18.}] (1) $k=\dfrac{\sqrt2}2$；(2) $d=\dfrac{2\sqrt2}3$；(3) 证明见解析
  \item[\textbf{19.}] (1) $b\geqslant2$；(2) $\dfrac13$
  \item[\textbf{20.}] (1) $a_5=8$；(2) $a_n=2n-1$
\end{itemize}

\gkpart{附加题}
\begin{itemize}[leftmargin=2.6em]
  \item[\textbf{21.}] A 证明见解析（$AB:AC=r_1:r_2$）；B $\boldsymbol\alpha=\begin{bmatrix}-1\\2\end{bmatrix}$；C $x-2y-4=0$；D $\left(-2,\dfrac43\right)$。
  \item[\textbf{22.}] (1) $AM=\dfrac{\sqrt{51}}5$；(2) $CM=\dfrac12$
  \item[\textbf{23.}] (1) $A_n=n-3$；(2) $B_n=\begin{cases}\dfrac{n(n-3)}6,&3\mid n,\\[4pt]\dfrac{(n-1)(n-2)}6,&3\nmid n.\end{cases}$
\end{itemize}
